\inputencoding{utf8}
\HeaderA{ghcnd\_stations}{GHCNd Station Data}{ghcnd.Rul.stations}
\keyword{datasets}{ghcnd\_stations}
%
\begin{Description}
GHCNd station data pulled from
\Rhref{https://www.ncei.noaa.gov/pub/data/ghcn/daily/ghcnd-stations.txt}{ghcnd-stations.txt}.
\end{Description}
%
\begin{Usage}
\begin{verbatim}
ghcnd_stations
\end{verbatim}
\end{Usage}
%
\begin{Format}
%
\begin{SubSection}{\code{ghcnd\_stations}}

A dataframe with 129658 rows and 9 columns:
\begin{description}

\item[ID] Station identification code
\item[LATITUDE] Station latitude in degrees
\item[LONGITUDE] Station longitude in degrees
\item[ELEVATION] Station elevation in meters
\item[STATE] U.S. postal code for the state (for U.S. stations only)
\item[NAME] Station name
\item[GSN\_FLAG] Flag that indicates whether the station is part of
the GCOS Surface Network (GSN). Unique values are blank and "GSN"
\item[HCN\_CRN\_FLAG] Flag that indicates whether the station is part of the
U.S. Historical Climatology Network (HCN) or U.S. Climate Reference
Network (CRN). Unique values are blank, "HCN", and "CRN"
\item[WMO\_ID] Station World Meteorological Organization (WMO) number

\end{description}

\end{SubSection}

\end{Format}
